
%(BEGIN_QUESTION)
% Copyright 2006, Tony R. Kuphaldt, released under the Creative Commons Attribution License (v 1.0)
% This means you may do almost anything with this work of mine, so long as you give me proper credit

Kromatografer virker ved å føre en væske eller gass gjennom en smal rørstrekning kalt en {\it kolonne}, og måle hva som kommer ut i andre enden. Denne strømmen består av en {\it bærefase} blandet med prøven som skal analyseres.

Forklar (mer detaljert) hva som skjer inne i "kolonnen" i en kromatograf, og hvordan dette henger sammen med analyse av kjemikalier i prøvestrømmen.

\underbar{file i00670no}
%(END_QUESTION)





%(BEGIN_ANSWER)

Når blandingen av bærefase og prøve beveger seg gjennom kolonnen, vil pakkmaterialet inne i kolonnen forsinke transporten av noen komponenter i prøven. Resultatet er at ulike kjemiske komponenter kommer ut av kolonnen på ulike tidspunkter. Grunnprinsippet i en kromatograf er at en flytende blanding av kjemikalier separeres over tid, slik at hver komponent kan måles kvantitativt med én og samme detektor.

\vskip 10pt

Kolonnen er vanligvis bare et langt metallrør med svært liten innvendig diameter. Noen ganger er forholdet mellom lengde og diameter enormt: ta for eksempel en pakket kolonne brukt for separasjon av komponenter i bensin, med kolonnelengde 100 meter og innvendig diameter bare 250 $\mu$m (mikrometer).

Som en sidebemerkning kommer dette eksemplet med en ekstremt lang kolonne fra Raymond P.W. Scotts utmerkede bok, {\it Principles and Practice of Chromatography}, tilgjengelig gratis på internett. I dette konkrete eksemplet med kromatografisk bensinanalyse var pakkmaterialet Petrocol DH med tykkelse 0.5 $\mu$m på innsiden av kolonneveggen, og prøvevolumet var 0.1 $\mu$l. Hele analysen tok omtrent 100 minutter, hovedsakelig på grunn av den lange kolonnen. Bæregassen var helium, og detektoren var av FID-typen (flammeionisering). Temperaturprogrammering ble brukt for å oppnå god separasjon av de tyngre komponentene, med start ved 35$^{o}$ C tidlig i analysen og slutt ved 200$^{o}$ C. Siden dette er gasskromatografi, gir høyere temperatur høyere gassviskositet (det motsatte av temperaturforholdet for væskeviskositet). Høyere temperatur mot slutten ga derfor større retardasjon av de tyngre gasskomponentene mot slutten av løpet, og dermed bedre separasjon.

%(END_ANSWER)





%(BEGIN_NOTES)


%INDEX% Måling, analytisk: kromatografi

%(END_NOTES)
